% =========================================================================
% 四、问题一：确定性日调度模型
% =========================================================================

\section{问题一：确定性日调度模型的建立与求解}

\subsection{决策变量与目标函数}

将一天离散为 144 个 10 分钟区间。决策变量为各区间的计划购电量 $g_t$、储能充电量 $c_t$、放电量 $d_t$、供能余量 $w_t$（$t=0,\dots,143$）与各区间起始储电量 $E_t$（$t=0,\dots,144$），共 $4\times144+145=721$ 个变量，全篇统一采用电量口径。以全天计划购电费用最小为目标：
\begin{equation}
  \min\ J=\sum_{t=0}^{143}p_t\,g_t .
\end{equation}

\subsection{约束条件}
\label{sec:q1-constraint}

\begin{align}
  &g_t+V_t+d_t=L_t+c_t+w_t, \qquad E_{t+1}=E_t+\eta_c c_t-\frac{d_t}{\eta_d}, \label{eq:q1-balance}\\
  &1200\le E_t\le 10800, \qquad 0\le c_t,\ d_t\le \frac{5000}{6}, \label{eq:q1-bound}\\
  &g_t\ge 0,\quad w_t\ge 0,\qquad E_0=E_{144}=6000. \label{eq:q1-cycle}
\end{align}
式~\eqref{eq:q1-balance} 第一式为区间能量平衡（供＝需＋余量），第二式为储能递推（充电按 $\eta_c$ 存入、放电按 $\eta_d$ 折损）；式~\eqref{eq:q1-bound} 依次为储电量允许范围（附录 1）与充放电功率上限（5000 kW 折算为区间电量上限 $5000/6$ kWh）；式~\eqref{eq:q1-cycle} 为题目要求的日循环条件（日初日末储电量相等，取附录 1 给定的 6000 kWh）。

\subsection{充放电互斥约束的松弛精确性}
\label{sec:q1-relax}

物理上要求 $c_t d_t=0$。以下命题保证无需引入二元变量。

\begin{quote}
  \textbf{命题.} 在 \S\ref{sec:q1-constraint} 的模型中，存在一个满足充放电互斥的最优解，即线性松弛在最优值意义下精确。
\end{quote}

\textbf{证明.} 设某最优解在区间 $t$ 同时充放电，取 $\delta=\min(c_t,\ d_t/0.81)>0$，将该区间充电减少 $\delta$、放电减少 $0.81\delta$。则储能变化量不变：
\begin{equation}
  -0.9\delta+\frac{0.81\delta}{0.9}=0,
\end{equation}
$E_t,E_{t+1}$ 及全部储能约束不受影响；能量平衡式中供给减少 $0.81\delta$、用能减少 $\delta$，净余量增加 $0.19\delta\ge0$，由 $w_t$ 吸收，可行性保持；$g_t$ 与目标值不变。重复此操作即得满足互斥且目标不劣的可行解。$\blacksquare$

该约化不依赖电价非负；其条件为余量可自由增加、效率乘积 $\eta_c\eta_d\le1$，以及模型存在有限最优解。数值上，求得的最优解同时充放电重叠量为 0，与命题一致。

\subsection{求解流程与数值校验}

模型为 721 变量、290 个等式约束的线性规划，调用 SciPy 的 \texttt{linprog}（HiGHS 求解器~\cite{highs,scipy}）求解，返回最优解，共 367 次迭代。求解流程为：① 读取附件 1 并按区间末端口径折算电量；② 以稀疏矩阵装配线性规划；③ 由 HiGHS 求解；④ 独立校验——能量平衡与储能递推的最大残差 $<10^{-12}$ kWh，储电量全程位于 $[1200,10800]$，$E_0=E_{144}=6000$，充放电功率无越界，同时充放电重叠量为 0，购电费用独立复算与求解器输出一致（差 $<10^{-10}$ 元）；⑤ 按官方模板写出 \texttt{result1.xlsx}。全部结果可由 \texttt{scripts/q1\_solve.py} 一条命令复现，核心程序见附录~\ref{app:q1-code}。

\subsection{结果分析}

\subsubsection{最优购电与储能调度结果}

\begin{table}[H]
  \centering
  \caption{指定时间段的计划购电量及全天结果（对应题目表 1）}
  \label{tab:q1-t1}
  \begin{tabular}{cccc}
    \toprule
    时间段 & 购电量 (kWh) & 时间段 & 购电量 (kWh) \\
    \midrule
    10:00--10:10 & 0.0000 & 16:00--16:10 & 445.4317 \\
    12:00--12:10 & 480.4124 & 18:00--18:10 & 531.8940 \\
    14:00--14:10 & 0.0000 & 20:00--20:10 & 0.0000 \\
    \midrule
    \multicolumn{2}{l}{全天购电量} & \multicolumn{2}{l}{\textbf{59482.6990 kWh}} \\
    \multicolumn{2}{l}{全天购电费} & \multicolumn{2}{l}{\textbf{35126.9486 元}} \\
    \bottomrule
  \end{tabular}
\end{table}

\begin{table}[H]
  \centering
  \caption{储能设备在指定时间段的充放电量及 0:00、24:00 储电量（对应题目表 2）}
  \label{tab:q1-t2}
  \begin{tabular}{ccc}
    \toprule
    时间段 & 充电量 (kWh) & 放电量 (kWh) \\
    \midrule
    0:00--4:00   & 4500.0000  & 0.0000 \\
    4:00--8:00   & 833.3333   & 6365.8412 \\
    8:00--12:00  & 4787.9643  & 1702.9970 \\
    12:00--16:00 & 5286.0352  & 91.1014 \\
    16:00--20:00 & 0.0000     & 5780.1319 \\
    20:00--24:00 & 5333.3333  & 2859.8681 \\
    \midrule
    0:00 储电量 & \multicolumn{2}{c}{6000.0000 kWh} \\
    24:00 储电量 & \multicolumn{2}{c}{6000.0000 kWh} \\
    \bottomrule
  \end{tabular}
\end{table}

最优策略（图~\ref{fig:q1-strategy}）呈清晰的``三充两放''节律：\textbf{凌晨低价段}以功率上限充电，并在全日最低价区间（5:30，0.3713 元/kWh）再次满功率充电，5:40 储电量首次触及上限 10800 kWh；\textbf{早高峰（5:50--9:20）}放电替代约 0.69--1.17 元/kWh 的购电，储电量回落至约 1860 kWh；\textbf{午间（约 9:30--15:10）}再次充电，能量来自午间相对低价购电（12--16 时购电 3810.0 kWh）与光伏富余吸收（同时段约 1384.9 kWh），14:30 前后再次触及上限；\textbf{晚高峰（18:00--20:40）}二次放电，覆盖电价最高的核心时段（峰值 1.3952 元/kWh），20:40 储电量触及下限 1200 kWh 并保持至 21:50；最后于\textbf{深夜低价段}满功率回充，24:00 恢复 6000 kWh 满足日循环约束。储电量两次触及上限、一次触及下限，说明容量约束在这些时刻处于活跃状态；但扩容是否经济，仍需通过敏感性分析定量判断（\S\ref{sec:q1-sens}）。

\begin{figure}[H]
  \centering
  \includegraphics[width=0.78\textwidth]{p1_fig2.pdf}
  \caption{最优购电与储能运行策略：计划购电量、储能充放电动作与储电量轨迹}
  \label{fig:q1-strategy}
\end{figure}

\subsubsection{与无储能基准的经济性比较}

无储能基准为逐区间即时平衡：$g_t^{0}=\max(L_t-V_t,0)$。

\begin{table}[H]
  \centering
  \caption{有无储能的经济性比较}
  \label{tab:q1-compare}
  \begin{tabular}{lccc}
    \toprule
    指标 & 无储能 & 有储能 & 变化 \\
    \midrule
    全天购电费 (元) & 48052.0466 & 35126.9486 & $-12925.0980$（节省率 26.90\%） \\
    全天购电量 (kWh) & 61789.9354 & 59482.6990 & $-2307.2364$（$-3.73\%$） \\
    \bottomrule
  \end{tabular}
\end{table}

能量机制与价格机制需分别解释。\textbf{能量机制}：储能充入 20740.6661 kWh、放出 16799.9396 kWh，往返损耗达 3940.7266 kWh；结合购电量减少 2307.2364 kWh，全天能量平衡意味着供能余量减少约 6247.96 kWh——储能吸收的富余供能超过自身损耗，购电量因此净下降。\textbf{价格机制}：按充放电量加权，放电替代电价约 1.1376 元/kWh、充电支付电价约 0.4930 元/kWh，比值约 2.31，大于往返效率决定的盈亏平衡比 $1/(\eta_c\eta_d)=1/0.81\approx1.23$，可见时段转移在存在能量损耗的条件下仍然盈利。

\subsubsection{容量与功率的敏感性分析}
\label{sec:q1-sens}

采用固定下限 1200 kWh、抬升储能上限的放宽路径（图~\ref{fig:q1-sens}a）：扩容右侧差分 $MV_+(h)=\left[C(E_{\max})-C(E_{\max}+h)\right]/h$ 在 $h=60,120,480$ kWh 时均为 0.6550 元/(kWh·日)，$h=1200$ kWh 时为 0.6539，价值函数在该区间接近线性：可用窗口从 9600 增至 10800 kWh，单日费用下降 784.68 元，区间平均扩容收益约 \textbf{0.6539 元/(kWh·日)}。功率上限从 5000 kW 翻倍至 10000 kW，费用仅下降 \textbf{0.47\%}（图~\ref{fig:q1-sens}b）。

\begin{figure}[H]
  \centering
  \includegraphics[width=0.88\textwidth]{p1_fig3.pdf}
  \caption{容量窗口与功率上限的敏感性（保持 $E_0=E_{144}=6000$ kWh）}
  \label{fig:q1-sens}
\end{figure}

\textbf{小结}：在本日数据与所考察的约束放宽路径下，容量窗口扩大的降费空间比继续提升充放电功率更为明显。需注意``元/kWh''与``元/kW''量纲不同，不能直接以数值大小比较两类投资的优先级；单日调度收益可为容量配置的讨论提供参考，但不足以单独证明投资合理性。
